\section{Problem definition and structural properties}\label{sec:prelim}

This section fixes the model, records its decision versions and the
parameterized objects we study, and collects the structural facts that
the later sections use as tools. The facts are elementary but used
throughout the subsequent proofs, so we prove them in full, with the
three longest proofs collected in Appendix~\ref{app:proofs}. Each is additionally checked
computationally, over the bounded instance families recorded in the
accompanying software archive.

An instance of MP-TSCFLP is given by four finite sets, the factories
$I$, the depots $J$, the customers $K$ and the products $L$, together
with nonnegative integer data. Opening factory $i$ costs $f_i$ and
opening depot $j$ costs $g_j$. Shipping one unit of product $l$ from
factory $i$ to depot $j$ costs $c_{ijl}$, and shipping it from depot $j$
to customer $k$ costs $d_{jkl}$. Factory $i$ can supply $b_{il}$ units of
product $l$ and depot $j$ can hold $p_{jl}$, while customer $k$ demands
$q_{kl}$. Both stages are complete bipartite, so no arc is forbidden a
priori. A solution opens a set of factories and depots, encoded by
$y\in\{0,1\}^{|I|}$ and $z\in\{0,1\}^{|J|}$, and routes the products through
nonnegative flows $x_{ijl}$ on the first stage and $w_{jkl}$ on the
second. We write $D_l=\sum_{k\in K}q_{kl}$ for the total demand of
product $l$ and $n=|I|+|J|$ for the number of candidate facilities.
Following \citet{MauriEtAl2021,PirkulJayaraman1998}, the problem
minimizes total cost,
\begin{align}
\min\;&\sum_{i\in I}f_iy_i+\sum_{j\in J}g_jz_j
      +\sum_{i,j,l}c_{ijl}\,x_{ijl}+\sum_{j,k,l}d_{jkl}\,w_{jkl}
\label{eq:obj}\\
\text{s.t.}\;
&\textstyle\sum_{j\in J}w_{jkl}\ge q_{kl}
   &&\forall k\in K,\ l\in L,\label{eq:C1}\\
&\textstyle\sum_{i\in I}x_{ijl}\ge\sum_{k\in K}w_{jkl}
   &&\forall j\in J,\ l\in L,\label{eq:C2}\\
&\textstyle\sum_{j\in J}x_{ijl}\le b_{il}\,y_i
   &&\forall i\in I,\ l\in L,\label{eq:C3}\\
&\textstyle\sum_{k\in K}w_{jkl}\le p_{jl}\,z_j
   &&\forall j\in J,\ l\in L,\label{eq:C4}\\
&y_i,z_j\in\{0,1\},\quad x_{ijl},w_{jkl}\ge 0.\label{eq:dom}
\end{align}
Constraint~\eqref{eq:C1} serves each demand, \eqref{eq:C2} balances flow
through a depot, and \eqref{eq:C3} and~\eqref{eq:C4} cap the throughput of
a factory and a depot by its capacity while forcing the facility open
before any flow may use it. At $y_i=0$ the right-hand side
of~\eqref{eq:C3} vanishes and nonnegativity forces $x_{ijl}=0$ for every
$j$ and $l$, and \eqref{eq:C4} acts on the $w_{jkl}$ in the same way at
$z_j=0$. Capacities and demands are per product and
never couple two products, a feature the analysis exploits throughout.
Figure~\ref{fig:network} shows the two-stage network and the decisions a
solution fixes.

% Requires: \usetikzlibrary{positioning,arrows.meta,shapes.geometric}
\begin{figure}[H]
  \centering
  \resizebox{0.86\textwidth}{!}{%
  \begin{tikzpicture}[
      >={Stealth[length=1.8mm]},
      thin,
      every node/.style={font=\footnotesize},
      fac/.style={draw=black!70,line width=0.5pt,fill=black!8,rounded corners=1pt,minimum size=6mm,inner sep=1pt},
      dep/.style={draw=black!70,line width=0.5pt,circle,fill=black!8,minimum size=6mm,inner sep=0pt},
      cus/.style={draw=black!70,line width=0.5pt,regular polygon,regular polygon sides=3,fill=black!4,minimum size=5mm,inner sep=0pt},
      lite/.style={draw=black!40,line width=0.35pt,->},
      p1/.style={draw=black,line width=0.7pt,->},
      p2/.style={draw=black,line width=0.7pt,dashed,->}
    ]
    % factories I
    \node[fac] (i1) at (0,2.1) {$1$};
    \node[fac] (i2) at (0,0.3) {$2$};
    % depots J
    \node[dep] (j1) at (3.3,3.0) {$1$};
    \node[dep] (j2) at (3.3,1.2) {$2$};
    \node[dep] (j3) at (3.3,-0.6) {$3$};
    % customers K
    \node[cus] (k1) at (6.6,3.0) {};
    \node[cus] (k2) at (6.6,1.2) {};
    \node[cus] (k3) at (6.6,-0.6) {};
    % layer captions
    \node[font=\scriptsize\itshape] at (0,3.3) {factories $I$};
    \node[font=\scriptsize\itshape] at (3.3,4.05) {depots $J$};
    \node[font=\scriptsize\itshape] at (6.6,4.05) {customers $K$};
    % stage-1 complete bipartite (light)
    \foreach \a in {i1,i2}{\foreach \b in {j1,j2,j3}{\draw[lite] (\a) -- (\b);}}
    % stage-2 complete bipartite (light)
    \foreach \a in {j1,j2,j3}{\foreach \b in {k1,k2,k3}{\draw[lite] (\a) -- (\b);}}
    % highlighted multiproduct arcs (two products)
    \draw[p1] (i1) -- (j2);
    \draw[p2] (i2) -- (j2);
    \draw[p1] (j2) -- (k2);
    \draw[p2] (j2) -- (k1);
    % annotations, each placed in a stroke-free pocket
    \node[anchor=south east,font=\scriptsize] at (i1.north west) {$y_1$};
    \node[anchor=north,font=\scriptsize] at (i2.south) {$b_{2l}$};
    \node[anchor=south,font=\scriptsize] at (j1.north) {$z_1$};
    \node[anchor=north,font=\scriptsize] at (j3.south) {$p_{3l}$};
    \node[font=\scriptsize] at (1.1,1.48) {$x_{ijl}$};
    \node[font=\scriptsize] at (3.98,1.88) {$w_{jkl}$};
    \node[anchor=west,font=\scriptsize] at (k2.east) {$q_{kl}$};
  \end{tikzpicture}%
  }
  \caption{The MP-TSCFLP network of
  model~\eqref{eq:obj}--\eqref{eq:dom}. A designer opens factories
  ($y_i$) and depots ($z_j$), after which each product flows across two
  capacitated stages, from factories through depots to customers, along
  stage-1 arcs carrying $x_{ijl}$ and stage-2 arcs carrying $w_{jkl}$.
  Factory capacities $b_{il}$ and depot capacities $p_{jl}$ bound the
  flow, which must meet every customer demand $q_{kl}$. The heavy solid
  arcs trace product $l=1$ and the heavy dashed arcs product $l=2$, so
  the two line styles convey the multiproduct structure, while the light
  arcs show that both stages are complete. The coupling of the two stages
  through the depots is the interaction that the reductions of
  Section~\ref{sec:complexity} turn into hardness.}
  \label{fig:network}
\end{figure}


A tiny instance already shows the interaction that drives the hardness.
Take one product and one factory of ample capacity, two depots and three
customers, with unit second-stage cost $d_{jk}$ equal to $0$ when depot
$j$ can reach customer $k$ cheaply and to $1$ otherwise, the product
index suppressed as there is a single product. Choosing which
depots to open so that every customer is served through a zero-cost arc
is exactly a set-cover decision, and no amount of first-stage freedom
removes it, because the single factory feeds whichever depots are open.
A core of the problem's combinatorial difficulty is thus already present
in the second stage, and the reductions of Section~\ref{sec:complexity}
grow from this observation.

The optimization problem has two decision versions. In the first only the
budget is bounded, and in the second the number of opened facilities is
bounded as well, which is the version we parameterize.

\begin{definition}[$\mathrm{MP\text{-}TSCFLP}(B)$]\label{def:B}
Given an MP-TSCFLP instance and an integer $B\ge 0$, decide whether some
feasible solution has total cost at most $B$.
\end{definition}

\begin{definition}[$\mathrm{MP\text{-}TSCFLP}(B,k)$]\label{def:Bk}
Given an MP-TSCFLP instance and integers $B\ge 0$ and $k\ge 0$, decide
whether some feasible solution has total cost at most $B$ and opens at
most $k$ facilities, that is $\sum_{i}y_i+\sum_{j}z_j\le k$. The
parameter is $k$.
\end{definition}

A tempting alternative would parameterize by the number of facilities
open in an optimal solution. Solution-dependent parameters are studied
in other settings, but this one is not well defined here, because the
optimum is not unique.
When a closed facility carries zero fixed charge, opening it yields
another optimum, since by Proposition~\ref{prop:mono} below the total
cost cannot rise and optimality of the original solution means it cannot
fall. On such instances the count can therefore
vary by up to $n$ depending on which optimum one names, so we do not
adopt it. Turning the count into an explicit
budget repairs the defect, and setting $k=n$ recovers
Definition~\ref{def:B}, so nothing is lost.

We phrase the complexity statements in the standard parameterized
vocabulary, which we recall briefly following
\citet{DowneyFellows2013} and \citet{Cygan2015}. A parameterized problem
is fixed-parameter tractable, in the class \FPT, when it is solvable in
time $h(p)\cdot\mathrm{poly}(N)$ for a computable function $h$ of the
parameter $p$, with $N$ the instance size, and it lies in \XP\ when the
polynomial degree may depend on $p$. A problem that is \NP-hard already at
a fixed value of $p$ is para-\NP-hard and, unless $\Ppoly=\NP$, admits
neither an \FPT\ nor an \XP\ algorithm. We call a parameterization
strongly para-\NP-hard when some fixed slice is strongly \NP-hard, so
that the exclusion survives a unary encoding.\footnote{The term is our
coinage for the combination, para-\NP-hardness witnessed by a strongly
\NP-hard slice, and we are not aware of an established name.} Hardness within the W
hierarchy is transferred by \FPT\ reductions. An \FPT\ reduction maps an
instance of one parameterized problem to an equivalent instance of
another in \FPT\ time, with the parameter of the image bounded by a
computable function of the source parameter
\citep{DowneyFellows2013,Cygan2015}. A parameterized problem is
W[2]-hard when some W[2]-complete problem, \textsc{Set Cover}
parameterized by the cover size for instance, admits an \FPT\ reduction
to it, and a W[2]-hard problem is fixed-parameter tractable only if
every problem in W[2] is, which is considered implausible. A kernel is a
polynomial-time self-reduction that maps an instance to an equivalent
instance of the same problem whose size is bounded by a computable
function of $p$ alone, and a polynomial kernel is one whose size bound is
polynomial. More generally, a polynomial compression is a polynomial-time
reduction that maps an instance of a parameterized problem to an
equivalent instance of some fixed problem, not necessarily the same one,
whose size is bounded by a polynomial of the parameter alone. A
polynomial kernel is thus the special case of a polynomial compression in
which the target problem is the problem itself.

The first structural fact isolates the routing subproblem. Once the open
facilities are fixed the remaining decisions are pure transport, and
because capacities and costs never couple products the transport splits
into one independent flow per product. Each such flow is an ordinary
minimum-cost flow on a layered network, which both makes the residual
value computable in strongly polynomial time and, by total unimodularity,
forces it to be an integer. A routing for the fixed design $(y,z)$ is a
pair of nonnegative flow vectors $(x,w)$ satisfying
\eqref{eq:C1}--\eqref{eq:C4} with $(y,z)$ held fixed, and its block $l$ is
the slice of $(x,w)$ formed by the variables of product $l$. We write
$v(y,z)$ for the least total cost of a routing for the design, a minimum
that Proposition~\ref{prop:oracle} shows to be attained whenever a routing
exists, with the convention $v(y,z)=+\infty$ when no routing exists, that
is when the design is infeasible.

\paragraph{The layered network} For a fixed design $(y,z)$ and a product
$l$, build the network $N_l(y,z)$ with a source $S$, a sink $T$, one node
per open factory ($y_i=1$), one node per customer, and, for each open
depot ($z_j=1$), a split pair $j_{\mathrm{in}}\!\to\! j_{\mathrm{out}}$.
The arcs are as follows. From the source, each $S\!\to\! i$ has capacity
$b_{il}$ and cost $0$. Each $i\!\to\! j_{\mathrm{in}}$ has infinite
capacity and cost $c_{ijl}$, the split arc
$j_{\mathrm{in}}\!\to\! j_{\mathrm{out}}$ has capacity $p_{jl}$ and cost
$0$, each $j_{\mathrm{out}}\!\to\! k$ has infinite capacity and cost
$d_{jkl}$, and each $k\!\to\! T$ has capacity $q_{kl}$ and cost $0$. A
closed facility contributes no node, so its flow is forced to $0$, in
accordance with the vanishing right-hand sides of \eqref{eq:C3}
and~\eqref{eq:C4} noted after the model. An open depot caps its throughput at $p_{jl}$
through the split arc, which is~\eqref{eq:C4} at $z_j=1$, and the
source-arc capacity $b_{il}$ of an open factory is~\eqref{eq:C3} at
$y_i=1$. Every $S$--$T$ path has the form
$S\!\to\! i\!\to\! j_{\mathrm{in}}\!\to\! j_{\mathrm{out}}\!\to\! k\!\to\!
T$, because every arc leads from one layer of nodes to the next and no
arc skips a layer or returns to an earlier one.
Figure~\ref{fig:oraclenet} draws $N_l(y,z)$ for a small design and
labels each arc family with its capacity, its cost and the model
variable it carries.

% Requires: \usetikzlibrary{positioning,arrows.meta,shapes.geometric}
\begin{figure}[H]
  \centering
  \resizebox{0.96\textwidth}{!}{%
  \begin{tikzpicture}[
      >={Stealth[length=1.8mm]},
      thin,
      every node/.style={font=\footnotesize},
      term/.style={draw=black!70,line width=0.5pt,circle,fill=black!15,minimum size=6mm,inner sep=1pt},
      fac/.style={draw=black!70,line width=0.5pt,fill=black!8,rounded corners=1pt,minimum size=6mm,inner sep=1pt},
      dep/.style={draw=black!70,line width=0.5pt,circle,fill=black!8,minimum size=7mm,inner sep=0pt},
      cus/.style={draw=black!70,line width=0.5pt,regular polygon,regular polygon sides=3,fill=black!4,minimum size=5mm,inner sep=0pt},
      capr/.style={draw=black!45,line width=0.35pt,->},
      cost/.style={draw=black,line width=0.7pt,->},
      splt/.style={draw=black,line width=0.9pt,->}
    ]
    % source and sink
    \node[term] (S) at (0,0.9) {$S$};
    \node[term] (T) at (11.6,0.9) {$T$};
    % open factory
    \node[fac] (i1) at (2.3,0.9) {$i$};
    % split pairs of the two open depots
    \node[dep] (j1in) at (4.7,2.3) {$1_{\mathrm{in}}$};
    \node[dep] (j1out) at (6.7,2.3) {$1_{\mathrm{out}}$};
    \node[dep] (j2in) at (4.7,-0.5) {$2_{\mathrm{in}}$};
    \node[dep] (j2out) at (6.7,-0.5) {$2_{\mathrm{out}}$};
    % customers
    \node[cus] (k1) at (9.2,2.3) {};
    \node[cus] (k2) at (9.2,-0.5) {};
    \node[font=\scriptsize,anchor=south] at (k1.north) {$k=1$};
    \node[font=\scriptsize,anchor=north] at (k2.south) {$k=2$};
    % layer captions
    \node[font=\scriptsize\itshape] at (0,3.4) {source};
    \node[font=\scriptsize\itshape] at (2.3,3.4) {open factory};
    \node[font=\scriptsize\itshape] at (5.7,3.4) {open depots, split pairs};
    \node[font=\scriptsize\itshape] at (9.2,3.4) {customers};
    \node[font=\scriptsize\itshape] at (11.6,3.4) {sink};
    % source arc: capacity b_il, cost 0
    \draw[capr] (S) -- (i1);
    \node[font=\scriptsize,anchor=south] at (1.15,1.02) {cap.\ $b_{il}$};
    % first-stage arcs: cost c_ijl, carry x_ijl; label inside the wedge
    \draw[cost] (i1) -- (j1in);
    \draw[cost] (i1) -- (j2in);
    \node[font=\scriptsize,align=center] at (3.85,0.9)
      {cost $c_{ijl}$\\ carries $x_{ijl}$};
    % split arcs: capacity p_jl, cost 0
    \draw[splt] (j1in) -- (j1out);
    \draw[splt] (j2in) -- (j2out);
    \node[font=\scriptsize,anchor=south] at (5.7,2.45) {cap.\ $p_{jl}$};
    \node[font=\scriptsize,anchor=north] at (5.7,-0.68) {cap.\ $p_{jl}$};
    % second-stage arcs: cost d_jkl, carry w_jkl; label below the network
    \draw[cost] (j1out) -- (k1);
    \draw[cost] (j1out) -- (k2);
    \draw[cost] (j2out) -- (k1);
    \draw[cost] (j2out) -- (k2);
    \node[font=\scriptsize,anchor=north] at (8.0,-1.05)
      {cost $d_{jkl}$, carries $w_{jkl}$};
    % demand arcs: capacity q_kl, cost 0; label above the upper arc
    \draw[capr] (k1) -- (T);
    \draw[capr] (k2) -- (T);
    \node[font=\scriptsize,anchor=south] at (10.55,1.97) {cap.\ $q_{kl}$};
  \end{tikzpicture}%
  }
  \caption{The layered network $N_l(y,z)$ of
  Proposition~\ref{prop:oracle} for one product $l$, drawn for a design
  with one open factory, two open depots and two customers. Source arcs
  $S\to i$ have capacity $b_{il}$ and cost $0$, each arc
  $i\to j_{\mathrm{in}}$ costs $c_{ijl}$ and carries the model variable
  $x_{ijl}$, each split arc $j_{\mathrm{in}}\to j_{\mathrm{out}}$ has
  capacity $p_{jl}$ and cost $0$, each arc $j_{\mathrm{out}}\to k$ costs
  $d_{jkl}$ and carries $w_{jkl}$, and each demand arc $k\to T$ has
  capacity $q_{kl}$ and cost $0$. This correspondence between the arc
  families and the model variables is the content of the proposition's
  proof.}
  \label{fig:oraclenet}
\end{figure}


\begin{proposition}[routing oracle]\label{prop:oracle}
Fix $(y,z)\in\{0,1\}^{|I|}\times\{0,1\}^{|J|}$. The residual routing problem
separates over products, and for each product $l$ its optimal cost equals
the minimum cost $\mu_l$ of an $S$--$T$ flow of value $D_l$ in the layered
network $N_l(y,z)$ defined above, with $\mu_l=+\infty$ when no such flow
exists. The finite $\mu_l$ are attained by integral flows and are
nonnegative integers, and $v(y,z)=\sum_l\mu_l$. Hence $v(y,z)$ is
computable in strongly polynomial time by $|L|$ minimum-cost flow
computations \citep{Orlin1993}.
\end{proposition}
The proof identifies each per-product block with a minimum-cost flow of
value $D_l$ in the layered network $N_l(y,z)$, matching routings and
flows in both directions at equal cost after a trimming step that
tightens \eqref{eq:C1} and~\eqref{eq:C2}. Total unimodularity of the
node--arc incidence matrix makes every finite block optimum an integer
attained by an integral flow. Summing the blocks delivers the identity
$v(y,z)=\sum_l\mu_l$ used throughout the paper. The full proof is in
Appendix~\ref{app:proofs}.

Whether a fixed design can serve the demand at all is settled without
solving any flow, by an aggregate test on each product. The completeness
of the two stages is what makes the test exact, since any open factory
reaches any open depot and any open depot reaches any customer, so only
the totals matter.

\begin{proposition}[feasibility]\label{prop:feas}
A design $(y,z)$ is feasible if and only if, for every product $l$,
\begin{equation}
\textstyle\sum_{i\in I}b_{il}\,y_i\ \ge\ D_l
\quad\text{(F1)}\qquad\text{and}\qquad
\sum_{j\in J}p_{jl}\,z_j\ \ge\ D_l\quad\text{(F2)}.
\label{eq:feas}
\end{equation}
The test runs in $O((n+|K|)\,|L|)$ time, of which $O(|K|\,|L|)$ computes
the totals $D_l$.
\end{proposition}
Necessity follows by aggregating the four constraint groups per
product, which bounds the served demand by the total open capacity of
each of the two layers, a two-cut argument. Sufficiency exhibits an
explicit flow of value $D_l$ in $N_l(y,z)$ through northwest-corner
transport constructions on the two complete stages, so the aggregate
test (F1) and (F2) is exact. The full proof is in Appendix~\ref{app:proofs}.

The characterization rests on the completeness of the two stages in an
essential way, and the following counterexample shows that
the dependence is real rather than an artefact of the proof.

\begin{remark}[completeness is essential]\label{rem:sparse}
The aggregate test fails once an arc may be forbidden, a forbidden arc
being modelled by deleting its flow variable. Take one product, two
factories, two depots and one customer, and open all four facilities.
Since there is a single product and a single customer we drop those
indices and write $b_i$, $p_j$ and $D$ for the capacities and the
demand, with $b_1=b_2=1$, $p_1=p_2=1$ and $D=2$, so (F1) and (F2) both
read $2\ge 2$ and hold. Suppose the first stage is sparse and both
factories may ship only to depot $1$. Depot $2$ then receives nothing,
so \eqref{eq:C2} forces it to forward nothing, while depot $1$ forwards
at most one unit by~\eqref{eq:C4}. The customer therefore receives at
most one unit, which violates \eqref{eq:C1} against the demand of two,
so no routing exists although both tests pass. The per-product totals no
longer determine feasibility, which is why the characterization is
stated only for the complete model.
\end{remark}

Opening more facilities preserves feasibility and never raises the
routing cost, only the fixed charge. This monotonicity underlies both the polynomial
cases below and the admissible bound used by the enumeration algorithm.

\begin{proposition}[monotonicity]\label{prop:mono}
If $(y,z)\le(y',z')$ componentwise then feasibility of $(y,z)$ implies
feasibility of $(y',z')$, and $v(y',z')\le v(y,z)$ always, with the
convention $v=+\infty$ on infeasible designs. In particular the all-open
design attains the least routing value among all feasible designs.
\end{proposition}
The feasibility part holds because the left-hand sides of (F1) and
(F2) are nondecreasing in the design. For the value part, $N_l(y,z)$ is
a subnetwork of $N_l(y',z')$, so every value-$D_l$ flow of the smaller
network is one of the larger network at the same cost and no block
optimum increases. Enlarging a design therefore preserves feasibility
and never raises the routing value, which is the form in which later
sections use the proposition. The full proof is in Appendix~\ref{app:proofs}.

Two regimes collapse to polynomial-time problems. When all fixed charges
vanish the fixed cost drops out and monotonicity makes opening everything
optimal, and when a single factory faces a single depot the routing is
forced and the optimum has a closed form.

\begin{proposition}[polynomial cases]\label{prop:fzero}
If $f\equiv 0$ and $g\equiv 0$, then MP-TSCFLP is solvable in strongly
polynomial time, its instance is feasible exactly when
$\sum_i b_{il}\ge D_l$ and $\sum_j p_{jl}\ge D_l$ for every $l$, and its
optimum is the all-open routing value $v(\mathbf 1,\mathbf 1)$.
\end{proposition}
\begin{proof}
With no fixed charge the cost of a design is its routing
value alone. By Proposition~\ref{prop:mono} the all-open design attains
the least routing value among feasible designs, and by
Proposition~\ref{prop:feas} it is feasible precisely under the displayed
conditions, which are (F1) and (F2) at $(y,z)=(\mathbf 1,\mathbf 1)$. If
they fail, no design at all is feasible, for every design lies
componentwise below $(\mathbf 1,\mathbf 1)$, so its feasibility would
imply that of the all-open design by the feasibility part of
Proposition~\ref{prop:mono}. In the feasible case the all-open design is
itself a solution of cost $v(\mathbf 1,\mathbf 1)$, and no feasible
design costs less, so the optimum equals $v(\mathbf 1,\mathbf 1)$,
computed in strongly polynomial time by Proposition~\ref{prop:oracle}.
\end{proof}

The second polynomial regime fixes one facility on each side. There the
design is forced and the optimum admits a closed form.

\begin{proposition}[single factory and depot]\label{prop:onebyone}
If $|I|=|J|=1$, then MP-TSCFLP is solvable in $O(|K|\,|L|)$ time after
reading the input, and feasibility reduces to $b_{1l}\ge D_l$ and
$p_{1l}\ge D_l$ for every $l$. When the instance is feasible and some
demand is positive the optimum equals
$f_1+g_1+\sum_{l}\bigl(c_{11l}D_l+\sum_k d_{1kl}q_{kl}\bigr)$, and when all
demand is zero the optimum is $0$, attained by opening nothing.
\end{proposition}
\begin{proof}
Feasibility forces the design, and nonnegativity of the costs then
makes the cheapest routing explicit. By
Proposition~\ref{prop:feas}, a design $(y_1,z_1)$
is feasible if and only if $b_{1l}y_1\ge D_l$ and $p_{1l}z_1\ge D_l$ for
every $l$, which is what (F1) and (F2) read with one facility per side.
If all demand is zero, these conditions hold for every design and the
stated condition holds because every $D_l$ equals $0$. In that case the
empty design is feasible, since (F1) and (F2) read $0\ge 0$ and the
all-zero routing satisfies \eqref{eq:C1}--\eqref{eq:C4}, and its cost $0$
is optimal by nonnegativity of the costs. Otherwise some $D_l>0$, and
any feasible design has $y_1=z_1=1$, since $y_1=0$ or $z_1=0$ would make
(F1) or (F2) read $0\ge D_l>0$. For such a design the conditions become
$b_{1l}\ge D_l$ and $p_{1l}\ge D_l$, so the instance is feasible exactly
under the stated condition, and feasibility forces the fixed charge to be
exactly $f_1+g_1$. The only remaining variables are $x_{11l}$ and
$w_{1kl}$. Feasibility needs $w_{1kl}\ge q_{kl}$ by~\eqref{eq:C1} and
$x_{11l}\ge\sum_k w_{1kl}\ge D_l$ by~\eqref{eq:C2}, and since all costs
are nonnegative the cheapest choice is $w_{1kl}=q_{kl}$ and
$x_{11l}=D_l$, of cost $\sum_k d_{1kl}q_{kl}$ on the second stage and
$c_{11l}D_l$ on the first. This choice is itself feasible, since
\eqref{eq:C3} reads $x_{11l}=D_l\le b_{1l}$ and \eqref{eq:C4} reads
$\sum_k w_{1kl}=D_l\le p_{1l}$, both guaranteed by the feasibility
condition, while \eqref{eq:C1} and~\eqref{eq:C2} hold with equality. No
feasible routing costs less, since its variables dominate this choice
componentwise by the bounds just derived and every cost coefficient is
nonnegative. Summing over $l$ and adding $f_1+g_1$ gives the closed
form, evaluated in $O(|K||L|)$ time.
\end{proof}

Finally, the cardinality version sits in \XP\ by brute force over the
opening pattern, a bound the algorithms of
Section~\ref{sec:algorithms} later refine and, under the Exponential
Time Hypothesis, match up to the constant in the exponent.

\begin{observation}[membership in \XP]\label{obs:xp}
$\mathrm{MP\text{-}TSCFLP}(B,k)$ is decidable in time
$n^{k+O(1)}\cdot\mathrm{poly}(N)$, with $N$ the instance size. When $n=0$ the empty
design is the only design, one application of
Propositions~\ref{prop:oracle} and~\ref{prop:feas} to it decides the
instance in polynomial time, and the counting below is not needed, so
assume $n\ge 1$. One may further normalize $k\le n$, since no design
opens more than $n$ facilities. Enumerating every set
$S\subseteq I\cup J$ of at most $k$ facilities then decides the
instance. Each $S$ is identified with the design $(y,z)$ whose
indicators it defines, so that $v(S)$ denotes $v(y,z)$, feasibility of
$S$ is tested by Proposition~\ref{prop:feas}, and for feasible $S$ the
total cost
$\sum_{i\in S\cap I}f_i+\sum_{j\in S\cap J}g_j+v(S)$ is evaluated with
$v(S)$ from
Proposition~\ref{prop:oracle}. The answer is \textsc{yes} if and only if
some such cost is at most $B$, since by the definition of $v$ the
cheapest solution whose opening pattern is $S$ costs exactly this
amount, and every solution opening at most $k$ facilities has its
pattern among the enumerated sets. The number of sets is
$\sum_{r=0}^{k}\binom{n}{r}\le\sum_{r=0}^{k}n^{r}\le(k+1)\,n^{k}$, using
$\binom{n}{r}\le n^{r}\le n^{k}$ for $r\le k$ and $n\ge 1$. Each set is
processed in polynomial time, and the factor $k+1\le n+1$ is absorbed
into the $n^{O(1)}$ term, which gives the stated bound.
\end{observation}
